%====================================================================
%   SECTION 4 — REGIME-AWARE MODELS
%====================================================================

\section{Regime-aware models}\label{sec:regimemodels}

We introduce the AR/HMM machinery on a generic series and then specialise to the spread $S_t$. The autoregressive model of order $p$ projects a series $Y_t$ on its $p$ most recent past values \parencite{hamilton1994},
\begin{equation}\label{eq:arp}
    Y_t \;=\; c \;+\; \phi_{1} Y_{t-1} \;+\; \cdots \;+\; \phi_{p} Y_{t-p} \;+\; \varepsilon_t, \qquad \varepsilon_t \overset{\mathrm{iid}}{\sim} \mathcal{N}(0,\sigma^{2}),
\end{equation}
with the simplest non-trivial case AR(1),
\begin{equation}\label{eq:ar1}
    Y_t \;=\; c + \phi\, Y_{t-1} + \varepsilon_t,
\end{equation}
having unconditional mean $c/(1-\phi)$ for $|\phi|<1$. AR(1) captures momentum ($\phi > 0$) or immediate mean reversion ($\phi < 0$) only if $(c,\phi,\sigma^{2})$ are constant over the sample; to relax that we let the parameters of \eqref{eq:ar1} be piecewise constant, switching with an unobserved discrete-time Markov chain $K_t \in \{1,\ldots,K\}$ \parencite{hamilton1989new}. The dynamics of $K_t$ are governed by a $K \times K$ transition matrix $\mathbf{P} = (p_{ij})$ and an initial distribution $\nu = (\nu_1,\ldots,\nu_K)^{\top}$,
\begin{equation}\label{eq:transitionmatrix}
    p_{ij} \;=\; \mathbb{P}(K_t = j \mid K_{t-1} = i),
    \qquad
    \nu_k \;=\; \mathbb{P}(K_1 = k),
\end{equation}
with the row-stochastic constraint $\sum_{j=1}^{K} p_{ij} = 1$ for every $i$. The Markov assumption is deliberately simple: the future depends on the past only through the current state, yet the framework is rich enough to capture persistence, sudden switching, and regime-dependent parameters within a single probabilistic model.

\subsection{Markov-switching AR(1) on the spread}\label{sec:msarspread}

We specialise to $S_t$ from \eqref{eq:rollingspreadret} and take $K = 2$. The joint generative model is
\begin{equation}\label{eq:msargen}
\begin{aligned}
    S_t \mid S_{t-1},\, K_t = k \;&\sim\; \mathcal{N}\!\bigl(c^{(k)} + \rho^{(k)} S_{t-1},\;(\sigma^{(k)})^{2}\bigr),\\[4pt]
    \mathbf{P} \;&=\; \begin{pmatrix} p_{11} & p_{12} \\ p_{21} & p_{22} \end{pmatrix},
\end{aligned}
\end{equation}
an AR(1) on the spread with all three parameters $(c^{(k)},\rho^{(k)},\sigma^{(k)})$ switching \parencite{krolzig1997msvar}. For $|\rho^{(k)}|<1$ the regime mean is $\mu^{(k)} = c^{(k)}/(1-\rho^{(k)})$, yielding the mean-deviation form
\begin{equation}\label{eq:meandevparam}
    S_t \;=\; \mu^{(k)} \;+\; \rho^{(k)}\bigl(S_{t-1} - \mu^{(k)}\bigr) \;+\; \varepsilon_t^{(k)},
    \qquad
    \varepsilon_t^{(k)} \;\sim\; \mathcal{N}\bigl(0,\,(\sigma^{(k)})^{2}\bigr).
\end{equation}
Values of $\rho^{(k)}$ close to zero correspond to fast mean reversion; values close to unity, to near-random-walk behaviour. Before estimation, $S_t$ is winsorised at $\mathrm{median}(S) \pm 4\cdot\mathrm{std}(S)$ and rescaled by $10^{4}$ for numerical stability.

\paragraph{Estimation, regime choice, and labelling.} The HMM parameter set
\begin{equation}\label{eq:thetadef}
    \Theta \;=\; \bigl\{\, \nu,\; \mathbf{P},\; c^{(k)},\; \rho^{(k)},\; \sigma^{(k)} : k=1,\ldots,K \,\bigr\}
\end{equation}
is estimated by maximum likelihood, fitted with the Markov-switching routine in the \texttt{statsmodels} package \parencite{seabold2010statsmodels} from multiple random starts to mitigate convergence to local maxima. The smoothed regime posteriors
\begin{equation}\label{eq:smoothedposterior}
    \gamma_t(k) \;=\; \mathbb{P}(K_t = k \mid S_1,\ldots,S_T;\Theta)
\end{equation}
fall out of the same fit. We fix $K = 2$ throughout, motivated by the regime-switching literature in which a parsimonious two-state model already captures the dominant non-stationarity in returns \parencite{hamilton1989new}; in preliminary exploration $K \in \{3, 4\}$ either failed to separate cleanly across seeds or split the high-volatility state into sub-regimes whose entry implications were indistinguishable at the relevant cost level. With $K=2$ we label regimes by innovation variance,
\begin{equation}\label{eq:regimelabels}
    k_{\mathrm{MR}} \;=\; \arg\min_{k}\,\sigma^{(k)},
    \qquad
    k_{\mathrm{DR}} \;=\; \arg\max_{k}\,\sigma^{(k)},
\end{equation}
giving a mean-reverting (MR) and a drifting (DR) regime, with $\gamma_t^{\mathrm{MR}} = \mathbb{P}(K_t = k_{\mathrm{MR}} \mid S_1,\ldots,S_T)$ and $\gamma_t^{\mathrm{DR}} = 1 - \gamma_t^{\mathrm{MR}}$. In live trading only past observations are available, so the filtered posterior $\mathbb{P}(K_t = k_{\mathrm{MR}} \mid S_1,\ldots,S_t)$ is used inside the walk-forward back-test while the smoothed posterior is reserved for diagnostic figures.

\subsection{Regime-aware strategies}\label{sec:regimestrategies}

We build two regime-aware strategies on top of the baseline \eqref{eq:baselinerule}. Both inherit the baseline's entry direction and exit logic without modification: at every bar the candidate position is $\pi_t^{\mathrm{baseline}}$ from \eqref{eq:baselinerule}, and the regime model can only veto or attenuate that candidate. Neither variant can generate a position the baseline did not generate, nor reverse a position the baseline took.

\paragraph{AR-gated.} A binary gate with danger threshold $\delta \in (0,1)$,
\begin{equation}\label{eq:argate}
    G_t \;=\; \mathbbm{1}\!\bigl\{\,\gamma_t^{\mathrm{MR}} \,\geq\, 1-\delta\,\bigr\},
\end{equation}
modifies the rule to
\begin{equation}\label{eq:arpositionrule}
    \pi_t \;=\;
    \begin{cases}
        0 & \text{if } G_t = 0 \quad\text{(liquidate, stay flat)},\\[2pt]
        \pi_t^{\mathrm{baseline}} & \text{if } G_t = 1.
    \end{cases}
\end{equation}
The gate is binary: trade exactly like the baseline, or sit in cash. It also applies to existing positions: if we are long the spread and the HMM flags danger, the position goes flat at the next bar.

\paragraph{MS-AR.} A continuous, regime-weighted entry threshold,
\begin{equation}\label{eq:msarthreshold}
    z_q^{\mathrm{eff}}(t) \;=\; \gamma_t^{\mathrm{MR}}\, z_q \;+\; \gamma_t^{\mathrm{DR}}\, z_v,
    \qquad\text{with } z_v \geq z_q,
\end{equation}
replaces $z_q$ in \eqref{eq:baselinerule} by $z_q^{\mathrm{eff}}(t)$. When MR is highly probable, $z_q^{\mathrm{eff}} \approx z_q$ and the strategy trades like the baseline; when DR dominates, a larger $|Z_t|$ deviation is required before entering.

Table~\ref{tab:strategy_definitions} summarises the four strategies compared throughout the rest of the thesis. All three active rules share the cointegration spread $S_t$ and its $z$-score $Z_t$ and differ only in whether and how $\gamma_t^{\mathrm{MR}}$ enters the decision. In particular, AR and MS-AR are not separate forecasters; on any bar where the HMM does not intervene they trade exactly like the baseline, so any difference in performance between baseline and the regime-aware variants is attributable entirely to the gate suppressing or attenuating baseline trades.

\begin{table}[H]
\centering
\caption{The four strategies compared in this thesis. The ``Gate'' column gives the indicator multiplying the baseline position $\pi_t^{\mathrm{baseline}}$ before it is taken; a value of $1$ means the baseline is always allowed to trade.}
\label{tab:strategy_definitions}
\small
\begin{tabular}{lllll}
\toprule
\textbf{Strategy} & \textbf{Always long?} & \textbf{Entry threshold} & \textbf{Gate} & \textbf{Reads HMM?} \\
\midrule
Buy \& Hold & Yes & n/a & n/a & No \\
Baseline    & No  & $z_q$ & $1$ & No \\
AR          & No  & $z_q$ & $\mathbbm{1}\{\gamma_t^{\mathrm{MR}} \geq 1-\delta\}$ & Yes (binary) \\
MS-AR       & No  & $\gamma_t^{\mathrm{MR}} z_q + \gamma_t^{\mathrm{DR}} z_v$ & $1$ & Yes (continuous) \\
\bottomrule
\end{tabular}
\end{table}
